 %%%%%%%%%%%%%%%---SYNOPSIS--%%%%%%%%%%%%%%%%%%%%%%%%%%%

%%%%%%%%%%%%%%%%%%%%%%%%%%%%%%%%%%%%%%%%%%%%%%%%%%%%%%
\documentclass[12pt,a4paper,oneside]{report}
%\documentclass[12pt,a4paper,twoside]{book}
\usepackage{chngcntr}
%\counterwithin{figure}{chapter}
%\usepackage[nottoc]{tocbibind}
\usepackage{changepage}
\pagestyle{plain}
\usepackage{makeidx}
\usepackage{nomencl}
\usepackage{titlesec}
\usepackage{epsf}
\usepackage{graphicx}
\usepackage{amsmath}
\usepackage{amsfonts}
\usepackage{epstopdf}
\usepackage{comment}
\usepackage[hidelinks]{hyperref}
\setcounter{secnumdepth}{3}
\setcounter{tocdepth}{3}
\usepackage{fancyhdr}
\usepackage[english]{babel}
\usepackage{blindtext}
\usepackage{upquote}
%\usepackage[titletoc]{appendix}
\usepackage{enumerate}
\usepackage{lipsum}
\usepackage[a4paper, inner=1.5cm, outer=3cm, top=2cm,bottom=3cm, bindingoffset=1cm]{geometry}
\textheight=235mm
\textwidth=145mm
\topmargin=0mm
\headheight=0mm
\headsep=0mm
\oddsidemargin=15mm
\evensidemargin=15mm
\makeindex
\makeglossary
\graphicspath{{./figures/}}
\usepackage{float}
\usepackage{microtype}
\usepackage[]{setspace}
\singlespacing
\usepackage{amssymb}
\setstretch{1}
\usepackage{cite}
\linespread{1.3}
\renewcommand{\thesection}{\arabic{section}}
\usepackage[english]{babel}
\addto{\captionsenglish}{%
  \renewcommand{\bibname}{\LARGE References}
}

%***********************************************************

\begin{document}
\newgeometry{top=1in,bottom=1in,right=1in,left=1in}
\thispagestyle{empty}
\begin{center}
    {\textbf{\textit{A project report on}\\}}
   {\Large\bf{COMPUTER SCIENCE PORTFOLIO WEBSITE} \\}

			    
			       \vspace{0.3in}
			     
			      
  {\textbf{ Submitted to:\\}}
		       
		       % \vspace{0.3 in}
               \vspace{0.4in}
			    {\large \bf GISMA UNIVERSITY OF APPLIED SCIENCES\\}

	%	       {\bf\textit{to be submitted by}}

			     \vspace{0.4in}
			    {\large \bf Samuel Ngugi Irungu \\
			    Reg. No. GH1055413\\}
			 
			  \vspace{0.3 in}

		       
		       
		       \vspace{0.3 in}

		      % {\bf\textit{of}}
		       
		     %  \vspace{0.3in}
			    {\Large{\bf Bachelor\\ of \\Computer Science}\\}
                   \vspace{0.3in}
               
			{\Large{\bf B201A COMPUTER SCIENCE LAB SSO326}\\}		    
                \vspace{0.3in}
	      \includegraphics[width=0.30\linewidth]{Gisma1.png}


			   %\vspace{0.2in}
		     
   
			 {\large\bf  Berlin Campus \\ 26 June, 2026}
			     
\end{center}
\date{}
\setcounter{page}{0}
\newpage
%****************************************************************
\section{Introduction:}
This report represents the development of a website created by HTML and CSS over the period of the quarter-semester.

\section{Project links}
\begin{enumerate}
    \item \bf{Portfolio Website Hosted on GitHub Webpages:} \href{https://samuelirungun.github.io/Portfolio/}{https://samuelirungu
    n.github.io/Portfolio/ \faLink}
    \item \bf{GitHub Repositories link:}\href{https://github.com/SamuelIrunguN/Portfolio}{https://github.com/SamuelIrunguN/Portfolio\faLink} 
    \item \bf{Python Projects:} \href{https://github.com/SamuelIrunguN/PythonProjects}{https://github.com/SamuelIrunguN/PythonProjects\faLink}
\end{enumerate}





% A promising  approach is the measurement of ultrasonic velocity
%**************************************************************
\section{Structure of Portfolio}
The portfolio is organized into 5 sections: Home, About, Projects, Skills and Contact.
In the "Home" section: the Hero section has the title Computer Science Student in bold and large text. It is followed with an introduction and personal motto below it. 
\noindent
The "about" section follows containing my interests and qualities.
The projects section contains three python projects: two geared towards Aerospace engineering projects and a miniature AI assistant that helps users to select a career based on weighted logic. 
\noindent
The "skills" section follows containing a list of skills followed by the "contact" section which is the last section, this is where the potential partners can contact me for project collaboration 



\section{Justification of Portfolio and CV Design}
The website was created based on a blend of primary colors: black and white, with a shade of orange from the horizons of the webpage. The black and white scheme was intended to match the CV; however, the same orange tint highlight was not easily readable on the CV, therefore I changed it to Olive Green.
\noindent
The colors black and white are to represent professionalism while the orange shades were to represent the "rise/dawning" of a new developer. In the Portfolio the orange colors emerge from the boundary/ horizons of the webpage to resemble the way the sun rises and sets. 

\noindent
The background is layered with a series of squares/tiles that blend in with the colors (dark grey and orange). These series of squares are intended to introduce the concept of vastness: the content and projects in my portfolio are geared towards computer science and aerospace engineering, therefore in combining the vastness of space and the systematic theme of coding I used a series of squares that span the whole portfolio.


%***************************************************************
\section{Tool Utilization}
I used VS Code as a code editor, in which I added HTML5, CSS3 AND Python 3.14 extensions. I used Git to add the files into my Git Hub repositories

%****************************************************************
%\section{Conclusions}

%****************************************************************

\section{Strengths and Weaknesses}
\subsection{Strengths}
\begin{enumerate}
    \item Consistent background-series of squares which makes it visually appealing
    \item Creative orange shade and tiles that blend together
    \item Clear categorized sections with consistent layout of white boxes containing text
    \item Projects are aligned to desired occupation i.e.
    
Working Computer Science Student:\\ ISAR Aerospace  \href{https://www.linkedin.com/jobs/view/4412623385/}{https://www.linkedin.com/jobs/view/4412623385/\faLink}


Working Computer Science Student Position at DLR Germany: \href{https://jobs.dlr.de/default/job/Student-%28mfd%29-Computer-Science%2C-Mathematics%2C-Data-Science%2C-Physics%2C-or-similar/5387-en_GB}{https://www.DLR Student position.com\faLink}
\end{enumerate}





\subsection{Weaknesses}
The design did not contain very appealing slide-ins
The olive-green from the CV breaks consistency from the portfolio website

%\pagebreak 
%****************************************************************

\section{Content and Features to add}
\subsection{Content}
\begin{enumerate}
 \item Replace current projects with more complex projects that are specific to my domain of interest.
 \end{enumerate}
 
\subsection{Features}
 \begin{enumerate}
     \item Add a slide-in hero card.
      \item Add a better hover effect for all the cards, that makes the unselected background blur
 \end{enumerate}
%%%%%%%%%%%%%%%%%%%%%%%%%%%%%%%%%%%%%%%%%%%%%%%%%
\section{Acknowledgment}
The author acknowledges the use of this template from "Synopsis of University Hyderabad" for writing the report
\end{document}

\end{document}
%*****************END*********